% Part: lambda-calculus
% Chapter: lambda-definability
% Section: fixpoints

\documentclass[../../../include/open-logic-section]{subfiles}

\begin{document}

\olfileid{lam}{ldf}{fp}

\olsection{نقاط ثابت}

فرض کنید بخواهیم تابع فاکتوریل را با بازگشت، به‌صورت ترمی چون
$\fn{Fac}$ با خاصیت زیر تعریف کنیم:
\[
\fn{Fac} \ident \lambd[n][\fn{IsZero}\, n\, \num 1 (\fn{Mult}\, n (\fn{Fac}(\fn{Pred} \, n)))]
\]
یعنی فاکتوریل~$n$ اگر~$n = 0$ باشد برابر~$1$، و در غیر این صورت برابر
$n$ ضرب‌در فاکتوریل~$n-1$ است. البته نمی‌توانیم ترم~$\fn{Fac}$ را به
این شیوه تعریف کنیم، زیرا خود~$\fn{Fac}$ در طرف راست رخ می‌دهد. چنین
تعریف‌های بازگشتیِ خودارجاع جزئی از حساب لامبدا نیستند. برای نمونه، در
تعریف ترم زیر
\[
\fn{Mult} \ident \lambd[ab][a (\fn{Add}\, b) \num 0]
\]
طرف راست فقط ترم‌هایی را در بر دارد که پیش‌تر تعریف شده‌اند، مانند
$\fn{Add}$. همواره می‌توانیم~$\fn{Add}$ را با ترم معرف آن جایگزین و
حاصل را به ترم لامبدای محضی برای~$\fn{Mult}$ تبدیل کنیم. اگر خود
$\fn{Add}$ ترم‌های تعریف‌شده‌ای را در بر داشته باشد (چنان‌که مثلاً
$\fn{Add}'$ دارد)، می‌توانیم این فرایند را ادامه دهیم تا سرانجام به یک
ترم لامبدای محض برسیم.

اما این امر دربارهٔ تعریف‌های بازگشتی مانند تعریف~$\fn{Fac}$ در بالا
درست نیست. اگر رخداد~$\fn{Fac}$ در طرف راست را با خود تعریف~$\fn{Fac}$
جایگزین کنیم، به عبارت زیر می‌رسیم:
\begin{align*}
  \fn{Fac} & \ident \lambd[n][\fn{IsZero}\, n\, \num{1}] \\
    & \qquad (\fn{Mult}\, n
  ((\lambd[n][\fn{IsZero} \, n \, \num 1\, (\fn{Mult}\, n\, (\fn{Fac}
    (\fn{Pred}\, n)))]) (\fn{Pred}\, n)))
\end{align*}
و هنوز~$\fn{Fac}$ را از طرف راست حذف نکرده‌ایم. روشن است که اگر این
فرایند را تکرار کنیم، تعریف پیوسته بلندتر می‌شود و هرگز ترم لامبدای
محضی به دست نمی‌آید. پس این شیوهٔ تعریف فاکتوریل، یا به‌طور کلی توابع
بازگشتی، عملی نیست.

بااین‌حال، تعریف بازگشتی نکته‌ای را به ما می‌گوید: اگر~$f$ ترمی معرف
تابع فاکتوریل باشد، آنگاه اعمال ترم
\[
\fn{Fac}' \ident \lambd[g][\lambd[n][\fn{IsZero} \, n \, \num 1\, (\fn{Mult} \, n\, (g (\fn{Pred} n)))]]
\]
بر~$f$، یعنی~$\fn{Fac}'\,f$، نیز تابع فاکتوریل را نمایش می‌دهد. یعنی
اگر~$\fn{Fac}'$ را تابعی بدانیم که یک تابع را می‌پذیرد و تابعی را
بازمی‌گرداند، مقدار~$\fn{Fac'}\, f$، مشروط بر اینکه~$f$ فاکتوریل باشد،
همان~$f$ است. تابع~$f$ با خاصیت~$\fn{Fac}' \, f \equal[\beta] f$ یک
\emph{نقطهٔ ثابت}~$\fn{Fac}'$ نامیده می‌شود. پس فاکتوریل نقطهٔ ثابت
$\fn{Fac}'$ است.

در حساب لامبدا ترم‌هایی وجود دارند که نقاط ثابت یک ترم داده‌شده را
محاسبه می‌کنند؛ سپس می‌توان از این ترم‌ها برای تبدیل ترمی مانند
$\fn{Fac}'$ به تعریف فاکتوریل استفاده کرد.

\begin{defn}
  \ollabel{defn:Turing-Y}
  \emph{ترکیبگر~Y} ترم زیر است:
  \[
  Y \ident (\lambd[ux][x(uux)])(\lambd[ux][x(uux)]).
  \]
\end{defn}

\begin{thm}
  برای هر ترم~$g$، ترم~$Y$ خاصیت~$Yg \red g(Yg)$ را دارد. بنابراین
  $Yg$ همواره نقطهٔ ثابت~$g$ است.
\end{thm}

\begin{proof}
  بگذارید~$U$ مخفف~$(\lambd[ux][x(uux)])$ باشد؛ پس~$Y \ident UU$.
  آنگاه
  \begin{align*}
    Y g &\ident (\lambd[ux][x(uux)])U\,g \\
    &\red (\lambd[x][x(UUx)])g\\
    & \red g(UUg) \ident g(Yg).
  \end{align*}
  چون~$g(Yg)$ و~$Yg$ هر دو به~$g(Yg)$ کاهش می‌یابند،
  $g(Yg) \equal[\beta] Yg$؛ پس~$Yg$ نقطهٔ ثابت~$g$ است.
\end{proof}

البته چون~$Yg$ یک سرخ‌عبارت است، کاهش می‌تواند بی‌پایان ادامه یابد:
\begin{align*}
  Y g &\red g (Y g) \\
    &\red g (g (Y g)) \\
    &\red g(g (g (Y g)))\\
    &\ldots
\end{align*}
پس می‌توان~$Yg$ را حاصل اعمال بی‌نهایت‌بارهٔ~$g$ بر خودش دانست. اگر
$g$ را یک بار دیگر بر آن اعمال کنیم، به‌اصطلاح، چیزی نمی‌افزاییم؛ اعمال
$g$ بر~$g$ که خود بی‌نهایت بار بر~$Yg$ اعمال شده، همچنان همان اعمال
بی‌نهایت‌بارهٔ~$g$ بر~$Yg$ است.

توجه کنید که دنبالهٔ بالا از گام‌های کاهش~$\beta$ که با~$Yg$ آغاز می‌شود
نامتناهی است. پس اگر~$Yg$ را بر ترمی اعمال کنیم، یعنی~$(Yg)N$ را در
نظر بگیریم، آن ترم نیز به بی‌نهایت ترم متفاوت، یعنی~$(g(Yg))N$،
$(g(g(Yg)))N$، \dots، کاهش می‌یابد. بااین‌همه، ممکن است دنبالهٔ
\emph{دیگری} از گام‌های کاهش در یک صورت نرمال پایان یابد.

برای مثال فاکتوریل را در نظر بگیرید. $\fn{Fac}$ را برابر
$Y\, \fn{Fac}'$ تعریف کنید (یعنی یک نقطهٔ ثابت~$\fn{Fac'}$). آنگاه:
\begin{align*}
  \fn{Fac}\,\num 3
  &\red Y\, \fn{Fac}' \, \num 3 \\
  &\red \fn{Fac}' (Y \,\fn{Fac}') \, \num 3 \\
  &\ident (\lambd[x][\lambd[n][\fn{IsZero} \, n\, \num 1\,
      (\fn{Mult}\, n\, (x (\fn{Pred}\, n)))]]) \, \fn{Fac} \, \num 3\\
  & \red \fn{IsZero} \, \num 3\, \num 1\,
      (\fn{Mult}\, \num 3\, (\fn{Fac} (\fn{Pred}\, \num 3))) \\
  &\red  \fn{Mult}\, \num 3 \, (\fn{Fac} \, \num 2).
  \intertext{به‌همین‌ترتیب،}
  \fn{Fac}\,\num 2
  &\red  \fn{Mult}\, \num 2 \, (\fn{Fac} \, \num 1) \\
  \fn{Fac}\,\num 1
  &\red  \fn{Mult}\, \num 1 \, (\fn{Fac} \, \num 0)
  \intertext{اما}
  \fn{Fac}\,\num 0
  &\red \fn{Fac}' (Y \,\fn{Fac}') \, \num 0 \\
  &\ident (\lambd[x][\lambd[n][\fn{IsZero} \, n\, \num 1\,
      (\fn{Mult}\, n\, (x (\fn{Pred}\, n)))]]) \, \fn{Fac} \, \num 0\\
  & \red \fn{IsZero} \, \num 0\, \num 1\,
      (\fn{Mult}\, \num 0\, (\fn{Fac} (\fn{Pred}\, \num 0))).\\
  &\red \num 1.
  \intertext{پس در مجموع}
  \fn{Fac}\,\num 3
  &\red  \fn{Mult}\, \num 3 \,
  (\fn{Mult} \, \num 2\, (\fn{Mult} \, \num 1\, \num 1)).
\end{align*}

آنچه برای~$\fn{Fac'}$ برقرار است برای هر تعریف بازگشتی نیز برقرار است.
فرض کنید معادلهٔ بازگشتی زیر را داریم:
\begin{align*}
g\,x_1\dots x_n & \equal[\beta] N
\intertext{که در آن~$N$ ممکن است~$g$ و~$x_1$، \dots،~$x_n$ را در بر
داشته باشد. آنگاه همواره ترمی چون
$G \ident (Y \lambd[g][\lambd[x_1\dots x_n][N]])$ وجود دارد، به‌طوری‌که}
G\,x_1 \dots x_n & \equal[\beta] \Subst{N}{G}{g}.
\intertext{زیرا بنا بر قضیهٔ نقطهٔ ثابت،}
G \ident (Y \lambd[g][\lambd[x_1\dots x_n][N]]) & \red \lambd[g][\lambd[x_1\dots x_n][N]](Y \lambd[g][\lambd[x_1\dots x_n][N]])\\
& \ident (\lambd[g][\lambd[x_1\dots x_n][N]])\,G
\intertext{و در نتیجه}
G\,x_1 \dots x_n 
& \red (\lambd[g][\lambd[x_1\dots x_n][N]])\,G\,x_1\dots x_n\\
& \red (\lambd[x_1\dots x_n][\Subst{N}{G}{g}])\,x_1\dots x_n\\
  & \red \Subst{N}{G}{g}.
\end{align*}

ترکیبگر~$Y$ در~\olref{defn:Turing-Y} از آلن تورینگ است. آلونزو چرچ
نسخهٔ دیگری پیشنهاد کرده بود که آن را~$Y_C$ می‌نامیم:
\[
Y_C \ident \lambd[g][(\lambd[x][g(xx)])(\lambd[x][g(xx)])].
\]
ترکیبگر چرچ اندکی ضعیف‌تر از ترکیبگر تورینگ است: داریم
$Y_Cg \equal[\beta] g(Y_Cg)$، اما رابطهٔ کاهش پیشروِ
$Y_Cg \bred g(Y_Cg)$ برقرار نیست؛ پس به‌ویژه کاهش یک‌گامی نیز در این
جهت نداریم. بگذارید~$V$ ترم
$\lambd[x][g(xx)]$ باشد؛ پس~$Y_Cg \ident (\lambd[g][VV])g$. آنگاه
\begin{align*}
  VV & \ident (\lambd[x][g(xx)])V \red g(VV)
  \text{ و بنابراین}\\
  Y_C g & \ident (\lambd[g][VV])g \red VV \red g(VV), \text{ اما همچنین}\\
  g(Y_C g) & \ident g((\lambd[g][VV])g) \red g(VV).
\end{align*}
به بیان دیگر، $Y_Cg$ و~$g(Y_Cg)$ به ترم مشترک~$g(VV)$ کاهش می‌یابند؛
پس~$Y_Cg \equal[\beta] g(Y_Cg)$. این خاصیت اغلب برای کاربردها کافی
است.

\end{document}
